% TOP_PROBLEM: {"id":30007531,"problem_number":"LOCAL-30007531","title":"Communication and forward guidance","category_id":21,"category":{"id":21,"name":"economics","display_name":"Economics","description":"Open economic theory statements, published research agendas, and proposed scoped specifications; question provenance and historical priority are qualified per record.","slug":"economics","order_index":21,"created_at":"2026-10-08T00:00:00Z"},"status":"research_question","external_url":null,"legacy_ids":[],"published":false,"statement":"Identify persistent communication effects on beliefs and actual spending, separating total effects from belief-mediated channels under explicit experimental assumptions.","economics":{"original_id":20,"field":"Inflation and monetary policy","kind":"identification","formalization_basis":"adapted_research_question","source_status":"research_agenda","priority_status":"earliest_found","priority_note":"This is the earliest explicit relevant statement verified in this audit, not a claim of original priority; older literature and earlier versions may contain antecedents.","earliest_verified_reference":{"authors":["Janet L. Yellen"],"title":"Macroeconomic Research After the Crisis","year":2016,"url":"https://www.federalreserve.gov/newsevents/speech/yellen20161014a.htm","doi":null,"locator":"Section “Inflation Dynamics,” paragraphs on whose expectations matter and how expectations form","evidence_summary":"The speech identifies unresolved questions about which agents’ inflation expectations matter, their quantitative influence, their formation, and the role of central-bank communication."},"current_reference":{"authors":["European Central Bank"],"title":"Research agenda: Monetary policy and its implementation","year":null,"url":"https://www.ecb.europa.eu/pub/economic-research/research_agenda/monetary_policy/html/index.en.html","doi":null,"locator":"“Forward guidance,” selected areas of focus","evidence_summary":"The agenda explicitly calls for theoretical and empirical analysis of forward-guidance effectiveness and optimal design."},"formalization_note":"The sources explicitly ask about expectation formation and forward-guidance effectiveness. The randomized design and mediation estimands are adapted specifications."},"statusQualification":"Published question or research agenda verified at the cited locator. The mathematical specification is adapted for this catalogue and includes additional assumptions; source publication does not establish first-ever priority or certify this exact model remains unsolved. The sources explicitly ask about expectation formation and forward-guidance effectiveness. The randomized design and mediation estimands are adapted specifications.","statusReviewedAt":"2026-10-08"}
% FIRST_POSTED: 2026-10-08
% ENTRY_KIND: statement_only
% !TeX program = lualatex
\documentclass[11pt]{article}
\usepackage{fontspec}
\usepackage[a4paper,margin=25mm]{geometry}
\usepackage{amsmath,amssymb,mathtools}
\usepackage{xurl}
\usepackage[hidelinks]{hyperref}
\newcommand{\sref}[2]{\hyperlink{source:#1:#2}{[S#2]}}
\setlength{\emergencystretch}{3em}
\begin{document}
% problemId:30007531
\section*{Communication and forward guidance}\label{30007531}
\subsection{Definitions and mathematical statement}
Fix a population of households or firms and a central-bank communication date. Randomize message \(Z_i\) among a control, a policy-path announcement, and an announcement with a credibility cue whose content is predeclared. Let \(B_i(z)\) be the resulting vector of beliefs about rates, inflation, and economic activity, and \(C_i(z)\) actual subsequent consumption or investment measured administratively. Define the total communication effect
\[
\tau_C(z)=E[C_i(z)-C_i(0)]
\]
and belief response \(\tau_B(z)=E[B_i(z)-B_i(0)]\). Identify these effects by randomized assignment, accounting for attrition and spillovers. To estimate a causal belief-to-action channel using assignment as an instrument, state monotonicity and exclusion restrictions: messages may convey central-bank information or alter confidence through pathways other than the measured beliefs.
Determine which message features produce persistent belief changes and actual behavior rather than immediate survey revisions. Compare effects across predetermined attention, liquidity, and trust measures, correcting for multiple subgroup comparisons. For extrapolation to aggregate forward guidance, specify equilibrium price feedback and interactions among treated agents; household-level treatment effects alone do not establish the aggregate response. The target is effective and credible communication in a declared environment, with explicit limits on transportability and mediation, rather than a universal forward-guidance coefficient.

\par\textbf{Formulation and scope.} The sources explicitly ask about expectation formation and forward-guidance effectiveness. The randomized design and mediation estimands are adapted specifications.

\par\textbf{Open-question provenance.} An explicit question or research agenda was verified at the source locator. This does not by itself certify that every later version remains unresolved \sref{30007531}{1}.

\par\textbf{Historical priority.} This is the earliest explicit relevant statement verified in this audit, not a claim of original priority; older literature and earlier versions may contain antecedents.

\subsection{Short English statement}
Identify persistent communication effects on beliefs and actual spending, separating total effects from belief-mediated channels under explicit experimental assumptions.

\subsection{Sources}
\begin{itemize}\raggedright
\item[\textnormal{[S1]}]\hypertarget{source:30007531:1}{} Janet L. Yellen. 2016. Macroeconomic Research After the Crisis. Section “Inflation Dynamics,” paragraphs on whose expectations matter and how expectations form.
\url{https://www.federalreserve.gov/newsevents/speech/yellen20161014a.htm}.
{\small\textit{Earliest verified question/agenda reference; historical priority qualified below. The speech identifies unresolved questions about which agents’ inflation expectations matter, their quantitative influence, their formation, and the role of central-bank communication.}}

\item[\textnormal{[S2]}]\hypertarget{source:30007531:2}{} European Central Bank. Year not verified. Research agenda: Monetary policy and its implementation. “Forward guidance,” selected areas of focus.
\url{https://www.ecb.europa.eu/pub/economic-research/research_agenda/monetary_policy/html/index.en.html}.
{\small\textit{Supporting or current literature; see evidence qualification. The agenda explicitly calls for theoretical and empirical analysis of forward-guidance effectiveness and optimal design.}}
\end{itemize}
\end{document}
